% ============================================================
%  RAPPORT DE SOUTENANCE — MOBIUS ROGUELIKE
%  M3G_STUDIO — Projet Jeu3D EPITA 2024-2025
% ============================================================
\documentclass[12pt, a4paper]{report}

% ── Encodage & Langue ──────────────────────────────────────
\usepackage[utf8]{inputenc}
\usepackage[T1]{fontenc}
\usepackage[french]{babel}

% ── Mise en page ───────────────────────────────────────────
\setlength{\headheight}{14pt}
\usepackage[top=2.5cm, bottom=2.5cm, left=2.8cm, right=2.8cm]{geometry}
\usepackage{setspace}
\onehalfspacing

% ── Graphiques & Couleurs ──────────────────────────────────
\usepackage{graphicx}
\usepackage{xcolor}
\usepackage{float}
\usepackage{caption}
\usepackage{subcaption}

% ── Tableaux ───────────────────────────────────────────────
\usepackage{booktabs}
\usepackage{array}
\usepackage{tabularx}
\usepackage{multirow}

% ── Listes & Code ──────────────────────────────────────────
\usepackage{enumitem}
\usepackage{listings}

% ── Liens & PDF ────────────────────────────────────────────
\usepackage[hidelinks]{hyperref}

% ── En-têtes / Pieds de page ───────────────────────────────
\usepackage{fancyhdr}
\usepackage{lastpage}

% ── Titres ─────────────────────────────────────────────────
\usepackage{titlesec}

% ── Divers ─────────────────────────────────────────────────

\usepackage{parskip}
\usepackage{mdframed}
\usepackage{amssymb}

% ╔══════════════════════════════════════════════════════════╗
%  COULEURS M3G_STUDIO
% ╚══════════════════════════════════════════════════════════╝
\definecolor{mobiuscyan}{RGB}{0,210,255}
\definecolor{mobiuspurple}{RGB}{100,40,160}
\definecolor{mobiusdark}{RGB}{15,15,35}
\definecolor{accentgold}{RGB}{255,200,0}
\definecolor{lightgray}{RGB}{245,245,248}
\definecolor{codesep}{RGB}{200,200,210}
\definecolor{okgreen}{RGB}{30,140,30}
\definecolor{warnorange}{RGB}{200,120,0}

% ╔══════════════════════════════════════════════════════════╗
%  STYLE DES TITRES
% ╚══════════════════════════════════════════════════════════╝
\titleformat{\chapter}[block]
  {\normalfont\huge\bfseries\color{mobiuspurple}}
  {\thechapter}{1em}{}
  [\vspace{2pt}\color{mobiuscyan}\hrule height 2pt]

\titleformat{\section}
  {\normalfont\Large\bfseries\color{mobiuspurple}}
  {\thesection}{0.8em}{}

\titleformat{\subsection}
  {\normalfont\large\bfseries\color{mobiusdark}}
  {\thesubsection}{0.6em}{}

% ╔══════════════════════════════════════════════════════════╗
%  EN-TÊTES / PIEDS DE PAGE
% ╚══════════════════════════════════════════════════════════╝
\pagestyle{fancy}
\fancyhf{}
\fancyhead[L]{\small\textcolor{mobiuspurple}{\textbf{MOBIUS Roguelike}} --- M3G\_STUDIO}
\fancyhead[R]{\small\textcolor{gray}{Rapport de soutenance}}
\fancyfoot[C]{\small\thepage\ / \pageref{LastPage}}
\renewcommand{\headrulewidth}{0.4pt}
\renewcommand{\footrulewidth}{0.2pt}

% ╔══════════════════════════════════════════════════════════╗
%  STYLE DE CODE
% ╚══════════════════════════════════════════════════════════╝
\lstset{
  language=Python,
  basicstyle=\ttfamily\small,
  keywordstyle=\color{mobiuspurple}\bfseries,
  stringstyle=\color{mobiuscyan},
  commentstyle=\color{gray}\itshape,
  backgroundcolor=\color{lightgray},
  frame=single,
  rulecolor=\color{codesep},
  breaklines=true,
  showstringspaces=false,
  numbers=left,
  numberstyle=\tiny\color{gray},
  numbersep=5pt,
}

% ╔══════════════════════════════════════════════════════════╗
%  BOÎTES INFO
% ╚══════════════════════════════════════════════════════════╝
\newmdenv[
  linecolor=mobiuscyan,
  linewidth=2pt,
  backgroundcolor=mobiuscyan!5,
  innerleftmargin=10pt,
  innerrightmargin=10pt,
  innertopmargin=8pt,
  innerbottommargin=8pt,
]{infobox}

% ============================================================
\begin{document}

% ╔══════════════════════════════════════════════════════════╗
%  PAGE DE TITRE
% ╚══════════════════════════════════════════════════════════╝
\begin{titlepage}
  \pagecolor{mobiusdark}
  \color{white}
  \centering
  \vspace*{1cm}

  
\includegraphics[width=0.55\textwidth]{mobius_logo_final.png}\\[1.2cm]

  {\fontsize{42}{50}\selectfont\bfseries\color{mobiuscyan} MOBIUS}\\[0.3cm]
  {\Large\bfseries\color{white} R O G U E L I K E}\\[0.6cm]
  \textcolor{accentgold}{\rule{10cm}{1.5pt}}\\[0.5cm]

  {\large\bfseries Rapport de soutenance --- Projet Jeu3D}\\[0.2cm]
  {\normalsize EPITA --- Promo 2027 --- 2024/2025}\\[1.2cm]

  \colorbox{mobiuspurple!80}{%
    \parbox{7cm}{\centering\vspace{4pt}
      {\large\bfseries M3G\_STUDIO}\\[2pt]
      {\small Studio de d\'{e}veloppement ind\'{e}pendant}
      \vspace{4pt}
    }
  }\\[1.5cm]

  \begin{tabular}{ll}
    \textbf{\color{mobiuscyan}Georges Boyer}       & Chef de groupe --- D\'{e}veloppement principal \\[4pt]
    \textbf{\color{mobiuscyan}Madeleine Diop}      & Environnement visuel \& graphismes \\[4pt]
    \textbf{\color{mobiuscyan}Emile Levast}        & Lore, narration \& site web \\[4pt]
    \textbf{\color{mobiuscyan}Ma\"{e}l Bosse-Plati\`{e}re} & Multijoueur en ligne \\
  \end{tabular}\\[2cm]

  \vfill
  \textcolor{gray}{\today}
\end{titlepage}
\nopagecolor
\color{black}

% ╔══════════════════════════════════════════════════════════╗
%  TABLE DES MATIÈRES
% ╚══════════════════════════════════════════════════════════╝
\tableofcontents
\newpage

% ============================================================
\chapter{Introduction}

Ce rapport de soutenance pr\'{e}sente l'avancement du projet \textbf{MOBIUS Roguelike}, r\'{e}alis\'{e} dans le cadre du projet \textit{Jeu3D} de premi\`{e}re ann\'{e}e \`{a} l'EPITA (promo 2027). Il retrace les choix techniques et cr\'{e}atifs effectu\'{e}s par l'\'{e}quipe M3G\_STUDIO, les fonctionnalit\'{e}s impl\'{e}ment\'{e}es, ainsi que le bilan organisationnel du groupe.

% ─────────────────────────────────────────────────────────────
\section{Pr\'{e}sentation du projet}

MOBIUS est un jeu de type \textit{roguelike} en vue du dessus (\textit{top-down shooter}) d\'{e}velopp\'{e} enti\`{e}rement en Python avec la biblioth\`{e}que Pygame. Le principe central est un voyage \`{a} travers six \'{e}poques historiques : le joueur affronte des vagues d'ennemis propres \`{a} chaque \`{e}re, r\'{e}cup\`{e}re des armes dans des coffres, et tente de survivre jusqu'\`{a} l'\`{e}re Futuriste.

Le jeu propose deux modes : un mode \textbf{solo} et un mode \textbf{multijoueur en ligne} (co-op 2 joueurs via UDP), chacun supportant la s\'{e}lection d'une classe de personnage parmi cinq disponibles.

\begin{infobox}
\textbf{Fiche technique}\\[4pt]
\begin{tabular}{ll}
  \textbf{Moteur / Langage} & Python 3.11+ / Pygame 2.x \\
  \textbf{R\'{e}solution}       & Plein \'{e}cran adaptatif (d\'{e}tection dynamique) \\
  \textbf{R\'{e}seau}           & UDP maison, port 55\,555, architecture client-serveur \\
  \textbf{\'{E}poques}          & 6 (Pr\'{e}histoire $\to$ Gr\`{e}ce $\to$ Edo $\to$ Moderne $\to$ WW2 $\to$ Futur) \\
  \textbf{Classes}          & 5 (Tank, Berserker, Vampire, Ninja, Archer) \\
  \textbf{Armes}            & Distanciel, m\^{e}l\'{e}e, AOE selon l'\'{e}poque \\
\end{tabular}
\end{infobox}

% ─────────────────────────────────────────────────────────────
\section{Pr\'{e}sentation du groupe}

Notre groupe M3G\_STUDIO est compos\'{e} de quatre \'{e}tudiants en premi\`{e}re ann\'{e}e \`{a} l'EPITA (promo 2027), r\'{e}unis autour d'une passion commune pour les jeux vid\'{e}o et la programmation.

\subsection{Georges Boyer --- Chef de groupe, d\'{e}veloppement principal}

Passionn\'{e} de programmation, Georges a naturellement pris en charge l'architecture g\'{e}n\'{e}rale du projet. Il a con\c{c}u le moteur de jeu (\texttt{base\_room.py}), l'orchestrateur principal (\texttt{main.py}), le syst\`{e}me de m\'{e}caniques (\texttt{mechanics.py}) et les constantes globales (\texttt{constants.py}). En tant que chef de groupe, il a \'{e}galement coordonn\'{e} les r\'{e}unions d'avancement et l'int\'{e}gration des diff\'{e}rents modules.

\textit{``Ce projet repr\'{e}sentait un vrai d\'{e}fi : passer de scripts Python isol\'{e}s \`{a} une architecture modulaire compl\`{e}te, capable de supporter un mode r\'{e}seau. C'est la partie qui m'a le plus appris sur la conception logicielle.''}

\subsection{Madeleine Diop --- Environnement visuel \& graphismes}

Madeleine a pris en charge tout le rendu visuel du jeu : cr\'{e}ation des fonds th\'{e}matiques par \'{e}poque, effets de particules (p\'{e}tales de cerisier dans l'\`{e}re Edo, pluie WW2, fum\'{e}e napol\'{e}onienne, scan holographique futuriste), et coh\'{e}rence graphique globale. Elle a d\'{e}velopp\'{e} le module \texttt{graphics.py} comprenant le \texttt{SpriteCache}, le \texttt{BackgroundRenderer} et les \texttt{ScreenEffects}.

\textit{``Rendre chaque \'{e}poque unique visuellement sans alourdir les performances a \'{e}t\'{e} mon d\'{e}fi principal. Les effets de particules diff\'{e}renci\'{e}s par \'{e}poque ont vraiment enrichi l'immersion.''}

\subsection{Emile Levast --- Lore, narration \& site web}

Emile s'est charg\'{e} de l'univers narratif du jeu et de la communication externe. Il a d\'{e}fini le \textit{lore} de MOBIUS (l'histoire du voyageur temporel, la logique des six \'{e}poques), r\'{e}dig\'{e} les textes du jeu, et d\'{e}velopp\'{e} le site web de M3G\_STUDIO.

\textit{``Construire un univers coh\'{e}rent autour d'un gameplay roguelike \'{e}tait un exercice d'\'{e}criture contraignant mais passionnant. Chaque \'{e}poque devait avoir sa propre logique narrative.''}

\subsection{Ma\"{e}l Bosse-Plati\`{e}re --- Multijoueur en ligne}

Ma\"{e}l a d\'{e}velopp\'{e} le module r\'{e}seau complet du jeu (\texttt{network.py}). Il a impl\'{e}ment\'{e} une architecture client-serveur UDP maison : le \texttt{GameServer} joue le r\^{o}le de serveur autoritaire, le \texttt{GameClient} envoie ses inputs et re\c{c}oit l'\'{e}tat du jeu pour l'afficher via le \texttt{ClientRenderer}.

\textit{``Impl\'{e}menter un protocole r\'{e}seau from scratch en UDP \'{e}tait de loin la partie la plus complexe. G\'{e}rer les timeouts, la d\'{e}synchronisation et la latence m'a donn\'{e} un aper\c{c}u concret des d\'{e}fis du multijoueur.''}

% ─────────────────────────────────────────────────────────────
\section{But du projet}

\subsection{Origine}

M3G\_STUDIO est n\'{e} de la volont\'{e} de cr\'{e}er un jeu original qui d\'{e}passe le simple exercice scolaire. L'id\'{e}e centrale --- traverser les grandes \'{e}poques de l'humanit\'{e} en affrontant des ennemis historiques --- est n\'{e}e d'une r\'{e}flexion sur ce qui rend un roguelike rejouable : la variation. En changeant radicalement d'univers visuel, d'armes et d'ennemis \`{a} chaque \'{e}poque, chaque run offre une exp\'{e}rience diff\'{e}rente.

\subsection{Nature du jeu}

MOBIUS est un \textit{top-down shooter roguelike} : le joueur \'{e}volue dans des ar\`{e}nes en vue a\'{e}rienne, \'{e}limine des vagues d'ennemis de difficult\'{e} croissante, et r\'{e}cup\`{e}re des am\'{e}liorations (power-ups) et armes dans des coffres. La m\'{e}canique de progression par \'{e}poques --- chacune avec ses propres ennemis, armes et ambiance visuelle --- est le fil conducteur du jeu.

\subsection{Objectif}

Livrer un jeu jouable, complet et soign\'{e} : menus anim\'{e}s, s\'{e}lection de classe, six \'{e}poques jouables, mode solo et multijoueur en ligne, et un rendu visuel coh\'{e}rent avec l'identit\'{e} de M3G\_STUDIO.

% ─────────────────────────────────────────────────────────────
\section{Pr\'{e}sentation de l'entreprise fictive}

\textbf{M3G\_STUDIO} a \'{e}t\'{e} fond\'{e} pour ce projet par les quatre membres de l'\'{e}quipe. Le nom reprend les initiales des fondateurs et affirme une identit\'{e} collective. Notre studio se sp\'{e}cialise dans les jeux ind\'{e}pendants \`{a} fort caract\`{e}re narratif et visuel, en particulier les roguelikes et les jeux d'action \`{a} progression.

Notre premier titre, \textbf{MOBIUS Roguelike}, illustre notre philosophie : un gameplay accessible mais profond, une direction artistique soign\'{e}e, et une rejouabilit\'{e} port\'{e}e par la variation des univers.

% ─────────────────────────────────────────────────────────────
\section{\'{E}tat de l'art \& inspirations}

\subsection{\'{E}tat de l'art}

Le genre roguelike top-down est largement repr\'{e}sent\'{e} dans l'industrie ind\'{e}pendante. Parmi les titres de r\'{e}f\'{e}rence :

\begin{itemize}[noitemsep]
  \item \textbf{Enter the Gungeon} (Dodge Roll, 2016) : r\'{e}f\'{e}rence du top-down shooter roguelike avec progression par armes.
  \item \textbf{Hades} (Supergiant Games, 2020) : roguelike narratif avec runs courts et progression m\'{e}ta.
  \item \textbf{Nuclear Throne} (Vlambeer, 2015) : roguelike action intense avec personnages aux capacit\'{e}s distinctes.
  \item \textbf{Vampire Survivors} (poncle, 2022) : popularisation du \textit{bullet heaven} avec progression automatique.
  \item \textbf{Call of Duty Zombies} (Treyarch, 2008) : inspiration pour le syst\`{e}me de vagues d'ennemis croissantes.
\end{itemize}

MOBIUS se diff\'{e}rencie par son syst\`{e}me d'\'{e}poques historiques qui transforme radicalement l'ambiance et les m\'{e}caniques \`{a} chaque niveau, et par son mode co-op en ligne int\'{e}gr\'{e}.

\subsection{Inspirations directes}

\textbf{Enter the Gungeon} a directement inspir\'{e} notre syst\`{e}me d'armes vari\'{e}es et la m\'{e}canique de coffres cachant des armes. La s\'{e}lection de classe avant le run s'inspire de \textbf{Hades}. La progression en vagues de difficult\'{e} croissante s'inspire de \textbf{Vampire Survivors} et de \textbf{Call of Duty Zombies}.

% ─────────────────────────────────────────────────────────────
\section{R\'{e}partition des t\^{a}ches}

\begin{table}[H]
\centering
\caption{R\'{e}partition des t\^{a}ches par membre}
\begin{tabularx}{\textwidth}{|l|X|c|}
\hline
\textbf{Membre} & \textbf{P\'{e}rim\`{e}tre} & \textbf{\'{E}tat} \\
\hline
Georges Boyer     & Architecture (base\_room, main, mechanics, constants), IA ennemis, syst\`{e}me d'armes, classes joueur, coffres, power-ups & \textcolor{okgreen}{\textbf{Termin\'{e}}} \\
\hline
Madeleine Diop    & Graphismes (graphics.py), effets par \'{e}poque, backgrounds, HUD, sprites & \textcolor{okgreen}{\textbf{Termin\'{e}}} \\
\hline
Emile Levast      & Lore des \'{e}poques, textes in-game, site web M3G\_STUDIO & \textcolor{okgreen}{\textbf{Termin\'{e}}} \\
\hline
Ma\"{e}l Bosse-Plati\`{e}re & Module r\'{e}seau UDP (network.py), synchronisation client-serveur, gestion d\'{e}connexion & \textcolor{warnorange}{\textbf{En cours}} \\
\hline
\end{tabularx}
\end{table}

% ─────────────────────────────────────────────────────────────
\section{Ressenti global du projet}

Le d\'{e}veloppement de MOBIUS a \'{e}t\'{e} une exp\'{e}rience intense et formatrice. Partir d'une id\'{e}e et aboutir \`{a} un jeu fonctionnel avec menu, multijoueur, six univers distincts et un syst\`{e}me de classes en quelques semaines nous a confront\'{e}s \`{a} des d\'{e}fis bien au-del\`{a} du simple cours de Python. Les difficult\'{e}s ont \'{e}t\'{e} nombreuses --- synchronisation r\'{e}seau, gestion des ressources visuelles, coh\'{e}rence des \'{e}tats du jeu --- mais chaque obstacle r\'{e}solu a renforc\'{e} notre ma\^{i}trise des outils et notre dynamique d'\'{e}quipe.

% ============================================================
\chapter{Avancement du projet}

% ─────────────────────────────────────────────────────────────
\section{Sc\'{e}nario et univers}

\subsection{Contexte narratif}

Le joueur incarne un voyageur temporel coinc\'{e} dans une boucle infinie --- une \textit{bande de M\"{o}bius} --- qui le fait traverser six \'{e}poques de l'histoire humaine. \`{A} chaque \'{e}poque, des cr\'{e}atures hostiles ont envahi la ligne temporelle et menacent l'int\'{e}grit\'{e} du continuum. Le joueur doit les \'{e}liminer pour passer \`{a} l'\'{e}poque suivante, en s'adaptant aux armes et aux techniques de combat propres \`{a} chaque \`{e}re.

\subsection{Progression par \'{e}poques}

\begin{table}[H]
\centering
\caption{\'{E}poque, th\`{e}me et multiplicateur de difficult\'{e}}
\begin{tabular}{clcl}
\toprule
\textbf{Ordre} & \textbf{\'{E}poque} & \textbf{Difficult\'{e}} & \textbf{Ambiance} \\
\midrule
1 & Pr\'{e}histoire       & $\times 1.0$ & \^{A}ge de pierre, cavernes, feux de camp \\
2 & Gr\`{e}ce Antique    & $\times 1.3$ & Temples, Olympe, soleil m\'{e}diterran\'{e}en \\
3 & P\'{e}riode Edo      & $\times 1.6$ & Japon m\'{e}di\'{e}val, cerisiers, nuit \'{e}toil\'{e}e \\
4 & \`{E}re Moderne      & $\times 2.0$ & Napol\'{e}on, champs de bataille, fum\'{e}e de canon \\
5 & Guerre Mondiale  & $\times 2.5$ & WW2, jungle, bunker, pluie \\
6 & \`{E}re Futuriste    & $\times 3.0$ & Station spatiale, n\'{e}ons, hologrammes \\
\bottomrule
\end{tabular}
\end{table}

\subsection{Transitions entre \'{e}poques}

Les transitions sont g\'{e}r\'{e}es par la classe \texttt{EpochTransition} dans \texttt{main.py} : un fondu noir progressif accompagn\'{e} du nom de l'\'{e}poque en grand texte color\'{e}. La progression des stats du joueur (PV, stamina, kills, coins) est conserv\'{e}e entre les \'{e}poques.

% ─────────────────────────────────────────────────────────────
\section{M\'{e}caniques de jeu}

\subsection{D\'{e}placements du joueur}

Le joueur se d\'{e}place au clavier (ZQSD) et vise \`{a} la souris. Un syst\`{e}me de \textbf{dash} (touche ESPACE) permet une esquive rapide avec cooldown (\texttt{DASH\_COOLDOWN = 30} frames) et co\^{u}t en stamina (\texttt{DASH\_STAMINA\_COST = 15}). La vitesse de base est modul\'{e}e par la classe choisie.

\subsection{Syst\`{e}me d'armes}

Les armes sont d\'{e}finies dans \texttt{WEAPONS\_DATA} et instanci\'{e}es via la classe \texttt{Weapon}. Chaque arme poss\`{e}de :
\begin{itemize}[noitemsep]
  \item Un type (\texttt{ranged} ou \texttt{melee})
  \item Des d\'{e}g\^{a}ts (\texttt{damage}), un cooldown (\texttt{cooldown\_max}) et un co\^{u}t en stamina
  \item Pour les armes \`{a} distance : vitesse de projectile (\texttt{projectile\_speed})
  \item Pour la m\^{e}l\'{e}e : port\'{e}e d'arc d'attaque (\texttt{melee\_range})
\end{itemize}

Les projectiles joueur (\texttt{Bullet}) h\'{e}ritent du sprite de l'arme, rotation\'{e} selon la direction de tir. Les attaques de m\^{e}l\'{e}e (\texttt{MeleeAttack}) g\'{e}n\`{e}rent une ellipse d'arc positionn\'{e}e dans la direction vis\'{e}e.

\subsection{Classes de personnage}

\begin{table}[H]
\centering
\caption{Classes disponibles et leurs caract\'{e}ristiques}
\begin{tabular}{lll}
\toprule
\textbf{Classe} & \textbf{Stats} & \textbf{Comp\'{e}tence sp\'{e}ciale} \\
\midrule
Tank       & 150 PV, Vitesse $-30\%$ & Bouclier : $-50\%$ d\'{e}g\^{a}ts re\c{c}us \\
Berserker  & 80 PV, Vitesse $+30\%$  & Rage : $\times 2$ d\'{e}g\^{a}ts pendant 5s \\
Vampire    & Stats normales           & $+10$ PV par kill (10s) \\
Ninja      & PV r\'{e}duits, Vitesse max & Invisibilit\'{e} br\`{e}ve \\
Archer     & Stats normales           & Fl\`{e}che multiple \`{a} distance \\
\bottomrule
\end{tabular}
\end{table}

\subsection{Power-ups et coffres}

Les \textbf{power-ups} (classe \texttt{PowerUp}) apparaissent al\'{e}atoirement et offrent quatre types de bonus : \textit{Damage}, \textit{Speed}, \textit{Health}, \textit{Stamina}. Ils pulsent visuellement et disparaissent apr\`{e}s 600 frames s'ils ne sont pas ramass\'{e}s.

Les \textbf{coffres} (classe \texttt{Chest}) contiennent une arme sp\'{e}cifique \`{a} l'\'{e}poque. Le joueur les ouvre en s'approchant (touche \texttt{E}), d\'{e}clenchant une animation d'ouverture et l'ajout de l'arme \`{a} son inventaire.

% ─────────────────────────────────────────────────────────────
\section{Fonctionnalit\'{e}s}

\subsection{Menu principal \& navigation}

Le menu principal (\texttt{MainMenuRenderer}) affiche le logo MOBIUS avec une animation flottante, des particules de fond et deux boutons : \textit{Jouer} et \textit{Quitter}. L'\'{e}cran de mode (\texttt{ModeSelectRenderer}) pr\'{e}sente deux grandes cartes cliquables --- \textbf{Solo} et \textbf{En ligne} --- avec effets de survol anim\'{e}s. Un slider de volume persistant (\texttt{VolumeSlider}) est pr\'{e}sent sur tous les \'{e}crans de menus.

\begin{figure}[H]
  \centering
  % Remplacer par : \includegraphics[width=0.75\textwidth]{screenshots/menu_principal.png}
  \fbox{\parbox{0.75\textwidth}{\centering\vspace{2.5cm}
    \textit{[Capture d'\'{e}cran : menu principal MOBIUS]}
  \vspace{2.5cm}}}
  \caption{\'{E}cran titre --- Menu principal MOBIUS}
\end{figure}

\subsection{S\'{e}lection de classe}

L'\'{e}cran \texttt{MenuRenderer} affiche cinq cartes de classe avec d\'{e}grad\'{e} de couleur, ic\^{o}ne, description et comp\'{e}tence sp\'{e}ciale. Un effet de surbrillance et d'\'{e}l\'{e}vation indique la carte survol\'{e}e. La s\'{e}lection se fait au clic ou via les touches \texttt{1}--\texttt{5}.

\subsection{Multijoueur en ligne}

D\'{e}taill\'{e} en section \ref{sec:multi}.

\subsection{HUD (affichage t\^{e}te haute)}

Le HUD affiche en temps r\'{e}el : barre de PV, barre de stamina, arme courante avec cooldown visuel, nom de l'\'{e}poque, num\'{e}ro de vague, compteur de kills et de pi\`{e}ces collect\'{e}es, et indicateur de comp\'{e}tence sp\'{e}ciale.

\begin{figure}[H]
  \centering
  % Remplacer par : \includegraphics[width=0.75\textwidth]{screenshots/hud_ingame.png}
  \fbox{\parbox{0.75\textwidth}{\centering\vspace{2.5cm}
    \textit{[Capture d'\'{e}cran : HUD en jeu --- \`{e}re Pr\'{e}historique]}
  \vspace{2.5cm}}}
  \caption{HUD pendant une partie}
\end{figure}

% ─────────────────────────────────────────────────────────────
\section{Multijoueur en ligne}
\label{sec:multi}

\subsection{Architecture r\'{e}seau}

Le module \texttt{network.py} impl\'{e}mente une architecture \textbf{client-serveur autoritaire sur UDP} (port 55\,555 par d\'{e}faut). Le \textbf{HOST} est le serveur : il simule l'int\'{e}gralit\'{e} du jeu et envoie l'\'{e}tat complet au client \`{a} chaque frame. Le \textbf{CLIENT} n'effectue aucune simulation locale : il envoie ses inputs (d\'{e}placement, tir, dash\ldots) et re\c{c}oit l'\'{e}tat du jeu via le \texttt{ClientRenderer}.

\subsection{Flux de connexion}

\begin{enumerate}[noitemsep]
  \item Le HOST lance le serveur UDP et attend le client (\'{e}cran \textit{En attente de P2\ldots}).
  \item Le CLIENT saisit l'IP du host (\texttt{IPInputRenderer}) et se connecte dans un thread s\'{e}par\'{e} pour ne pas bloquer l'UI.
  \item Une fois connect\'{e}s, chaque joueur arrive sur l'\'{e}cran de s\'{e}lection de classe.
  \item Le HOST envoie le signal \texttt{send\_start} ; le jeu d\'{e}marre simultan\'{e}ment.
\end{enumerate}

\subsection{Synchronisation}

\'{E}l\'{e}ments synchronis\'{e}s \`{a} chaque frame :
\begin{itemize}[noitemsep]
  \item Position et direction du joueur P2
  \item Points de vie de P2
  \item \'{E}tat des ennemis (position, PV)
  \item Projectiles ennemis actifs
  \item Transitions d'\'{e}poques et \'{e}v\'{e}nements game-over
\end{itemize}

La gestion des timeouts (8 secondes) permet de d\'{e}tecter proprement une d\'{e}connexion.

\subsection{D\'{e}fis rencontr\'{e}s}

L'impl\'{e}mentation r\'{e}seau a pos\'{e} plusieurs probl\`{e}mes techniques : la s\'{e}rialisation efficace de l'\'{e}tat du jeu (nombreux sprites), la gestion de la latence variable en UDP, et la synchronisation des transitions d'\'{e}poque entre les deux machines. La solution retenue --- serveur autoritaire sans pr\'{e}diction client --- privil\'{e}gie la simplicit\'{e} et la coh\'{e}rence.

% ─────────────────────────────────────────────────────────────
\section{Graphismes et environnement visuel}

\subsection{D\'{e}cors par \'{e}poque}

Chaque \'{e}poque dispose d'une classe h\'{e}ritant de \texttt{BaseRoom} et impl\'{e}mentant \texttt{draw\_epoch\_decoration()} :

\begin{table}[H]
\centering
\caption{Effets visuels par \'{e}poque}
\begin{tabular}{lll}
\toprule
\textbf{\'{E}poque} & \textbf{Couleur th\'{e}matique} & \textbf{Effet de d\'{e}coration} \\
\midrule
Pr\'{e}histoire   & Brun (139, 69, 19)     & Flammes anim\'{e}es (optionnelles) \\
Gr\`{e}ce Antique & Or grec (200, 180, 80) & Soleil anim\'{e} (optionnel) \\
P\'{e}riode Edo   & Rouge (180, 30, 30)    & P\'{e}tales de cerisier flottants \\
\`{E}re Moderne   & Bleu (0, 80, 160)      & Volutes de fum\'{e}e de canon \\
Guerre Mondiale & Vert (60, 80, 40)      & Pluie battante \\
\`{E}re Futuriste & Cyan (0, 220, 255)     & Scan holographique + chiffres 0/1 \\
\bottomrule
\end{tabular}
\end{table}

\subsection{Syst\`{e}me de sprites et cache}

Le module \texttt{graphics.py} centralise le chargement des assets via la classe \texttt{SpriteCache} (singleton) qui \'{e}vite les rechargements redondants. Le \texttt{BackgroundRenderer} g\'{e}n\`{e}re les fonds th\'{e}matiques par \'{e}poque. Les \texttt{ScreenEffects} g\`{e}rent les flashs d'impact, les fades et les effets de transition.

\begin{figure}[H]
  \centering
  \begin{subfigure}[b]{0.47\textwidth}
    \fbox{\parbox{\textwidth}{\centering\vspace{2cm}
      \textit{[Capture : \`{e}re Edo --- p\'{e}tales de cerisier]}
    \vspace{2cm}}}
    \caption{\`{E}re Edo}
  \end{subfigure}
  \hfill
  \begin{subfigure}[b]{0.47\textwidth}
    \fbox{\parbox{\textwidth}{\centering\vspace{2cm}
      \textit{[Capture : \`{e}re Futuriste --- scan holographique]}
    \vspace{2cm}}}
    \caption{\`{E}re Futuriste}
  \end{subfigure}
  \caption{Exemples de d\'{e}corations d'\'{e}poque}
\end{figure}

% ─────────────────────────────────────────────────────────────
\section{Architecture du code}

\subsection{Structure modulaire}

\begin{verbatim}
mobius/
├── main.py               # Orchestrateur, états, menus, boucle principale
├── core/
│   ├── constants.py      # Couleurs, données armes, skills, epochs
│   ├── graphics.py       # SpriteCache, BackgroundRenderer, ScreenEffects
│   ├── mechanics.py      # Weapon, Bullet, MeleeAttack, PowerUp, Chest
│   ├── base_room.py      # Classe mère de toutes les salles
│   └── network.py        # GameServer, GameClient, ClientRenderer
└── epoques/
    ├── prehistoire.py    # Salle ère Préhistorique
    ├── grece.py          # Salle Grèce Antique
    ├── edo.py            # Salle Période Edo
    ├── moderne.py        # Salle Ère Moderne
    ├── contemporain.py   # Salle Guerre Mondiale
    └── futuristique.py   # Salle Ère Futuriste
\end{verbatim}

\subsection{Machine \`{a} \'{e}tats}

La classe \texttt{Game} g\`{e}re l'\'{e}tat global du jeu via des constantes enti\`{e}res :

\begin{table}[H]
\centering
\caption{\'{E}tats du jeu et transitions}
\begin{tabular}{lll}
\toprule
\textbf{\'{E}tat} & \textbf{Description} & \textbf{Transition vers} \\
\midrule
\texttt{MAIN\_MENU}        & \'{E}cran titre            & \texttt{MODE\_SELECT} \\
\texttt{MODE\_SELECT}      & Solo ou En ligne       & \texttt{CHARACTER\_SELECT} / \texttt{ONLINE\_MENU} \\
\texttt{ONLINE\_MENU}      & H\'{e}berger / Rejoindre   & \texttt{WAITING\_CLIENT} / \texttt{IP\_INPUT} \\
\texttt{CHARACTER\_SELECT} & Choix de classe        & \texttt{PLAYING} \\
\texttt{PLAYING}           & Jeu en cours           & \texttt{GAME\_OVER} \\
\texttt{GAME\_OVER}        & Fin de partie          & \texttt{MAIN\_MENU} / \texttt{PLAYING} \\
\bottomrule
\end{tabular}
\end{table}

\subsection{Extrait de code repr\'{e}sentatif}

\begin{lstlisting}[caption=M\'{e}thode change\_epoch() dans main.py --- transition d'\'{e}poque avec persistance des stats]
def change_epoch(self, next_epoch: str | None):
    if next_epoch and next_epoch in self.rooms:
        p = self.current_room.player
        stats = {"skill": p.skill, "kills": p.kills,
                 "coins": p.coins, "health": p.health,
                 "max_health": p.max_health}
        def _do():
            self.current_epoch = next_epoch
            self.current_room  = self.rooms[next_epoch]
            self.current_room.start(
                stats["skill"], mode=self.net_mode,
                player_stats=stats, server=self.server)
            if self.server:
                self.server.send_start(stats["skill"])
        self.transition.start(next_epoch, _do)
    else:
        self.game_state = GAME_OVER
\end{lstlisting}

% ============================================================
\chapter{Bilan de l'avancement}

\section{Ce qui a \'{e}t\'{e} r\'{e}alis\'{e}}

L'ensemble des fonctionnalit\'{e}s pr\'{e}vues en priorit\'{e} haute a \'{e}t\'{e} livr\'{e} :

\begin{itemize}
  \item[$\checkmark$] Six \'{e}poques jouables avec d\'{e}cors distincts, ennemis et armes th\'{e}matiques
  \item[$\checkmark$] Cinq classes de personnage avec comp\'{e}tences sp\'{e}ciales
  \item[$\checkmark$] Syst\`{e}me d'armes complet (distanciel, m\^{e}l\'{e}e, projectiles ennemis)
  \item[$\checkmark$] Coffres et power-ups avec animations
  \item[$\checkmark$] Menus anim\'{e}s complets (titre, mode, classe, game over)
  \item[$\checkmark$] Transitions d'\'{e}poques avec fondu et texte anim\'{e}
  \item[$\checkmark$] Mode solo pleinement fonctionnel
  \item[$\checkmark$] Module r\'{e}seau UDP avec architecture client-serveur
  \item[$\checkmark$] HUD complet (PV, stamina, arme, vague, kills, coins)
  \item[$\checkmark$] Logo et identit\'{e} visuelle M3G\_STUDIO (pixel-art cyan/violet)
  \item[$\checkmark$] Site web M3G\_STUDIO (Emile Levast)
\end{itemize}

\section{Difficult\'{e}s et ajustements}

\begin{itemize}
  \item \textbf{Flammes pr\'{e}historiques et soleil grec} : ces effets ont \'{e}t\'{e} d\'{e}sactiv\'{e}s pour raisons de performance. Ils restent comment\'{e}s dans le code et sont r\'{e}activables facilement.
  \item \textbf{Synchronisation r\'{e}seau compl\`{e}te} : le rendu client est fonctionnel mais certains effets de particules ne sont pas transmis. La synchronisation des sons reste partielle.
  \item \textbf{Mode LAN vs Internet} : le multijoueur fonctionne de fa\c{c}on fiable en r\'{e}seau local. Sur Internet, l'ouverture du port UDP 55\,555 reste n\'{e}cessaire c\^{o}t\'{e} h\^{o}te.
\end{itemize}

\section{\'{E}carts pr\'{e}visionnel / r\'{e}el}

\begin{table}[H]
\centering
\caption{Comparaison pr\'{e}visionnel vs r\'{e}alis\'{e}}
\begin{tabularx}{\textwidth}{|X|c|c|}
\hline
\textbf{Fonctionnalit\'{e}} & \textbf{Pr\'{e}vu} & \textbf{R\'{e}alis\'{e}} \\
\hline
6 \'{e}poques jouables              & Oui & \textcolor{okgreen}{Oui} \\
\hline
5 classes joueur                & Oui & \textcolor{okgreen}{Oui} \\
\hline
Mode solo                       & Oui & \textcolor{okgreen}{Oui} \\
\hline
Mode multijoueur LAN            & Oui & \textcolor{okgreen}{Oui} \\
\hline
Mode multijoueur Internet       & Oui & \textcolor{warnorange}{Partiel (port forwarding requis)} \\
\hline
Effets d\'{e}cor toutes \'{e}poques    & Oui & \textcolor{warnorange}{4/6 actifs (Pr\'{e}histoire/Gr\`{e}ce d\'{e}sactiv\'{e}s)} \\
\hline
Site web M3G\_STUDIO            & Oui & \textcolor{okgreen}{Oui} \\
\hline
Sons / Musique                  & Bonus & \textcolor{warnorange}{Partiellement int\'{e}gr\'{e}} \\
\hline
\end{tabularx}
\end{table}

% ============================================================
\chapter{Communication}

\section{Site web M3G\_STUDIO}

Le site web de M3G\_STUDIO a \'{e}t\'{e} r\'{e}alis\'{e} par Emile Levast. Il pr\'{e}sente le studio, le jeu MOBIUS, les membres de l'\'{e}quipe, et propose un lien de t\'{e}l\'{e}chargement du jeu. Le site refl\`{e}te l'identit\'{e} visuelle cyan/violet de MOBIUS et inclut des pages d\'{e}di\'{e}es au lore des six \'{e}poques et aux contr\^{o}les du jeu.

\begin{figure}[H]
  \centering
  % Remplacer par : \includegraphics[width=0.75\textwidth]{screenshots/site_web.png}
  \fbox{\parbox{0.75\textwidth}{\centering\vspace{2.5cm}
    \textit{[Capture d'\'{e}cran : page d'accueil du site M3G\_STUDIO]}
  \vspace{2.5cm}}}
  \caption{Site web M3G\_STUDIO}
\end{figure}

\section{Logo et identit\'{e} visuelle}

Le logo MOBIUS a \'{e}t\'{e} con\c{c}u en pixel-art avec un d\'{e}grad\'{e} cyan vers bleu profond, agr\'{e}ment\'{e} d'un effet de foudre rose en bas. Ce logo est utilis\'{e} de fa\c{c}on coh\'{e}rente sur l'\'{e}cran titre, dans tous les menus, sur le site web et dans les documents du projet.

\begin{figure}[H]
  \centering
  
\includegraphics[width=0.45\textwidth]{mobius_logo_final.png}
  \caption{Logo officiel MOBIUS --- M3G\_STUDIO}
\end{figure}

\section{Manuel d'utilisation}

Les contr\^{o}les du jeu sont affich\'{e}s directement dans l'interface de s\'{e}lection de mode, garantissant que le joueur est inform\'{e} avant de commencer.

\begin{table}[H]
\centering
\caption{Contr\^{o}les du jeu}
\begin{tabular}{ll}
\toprule
\textbf{Touche} & \textbf{Action} \\
\midrule
\texttt{ZQSD}        & D\'{e}placement \\
\texttt{Clic gauche} & Tirer / Attaquer \\
\texttt{ESPACE}      & Dash (esquive) \\
\texttt{F}           & Comp\'{e}tence sp\'{e}ciale \\
\texttt{E}           & Ouvrir un coffre \\
\texttt{1 / 2}       & Changer d'arme \\
\texttt{ESC}         & Menu pause / Retour \\
\bottomrule
\end{tabular}
\end{table}

% ============================================================
\chapter{Conclusion}

Le d\'{e}veloppement de MOBIUS Roguelike a repr\'{e}sent\'{e} pour M3G\_STUDIO bien plus qu'un exercice acad\'{e}mique. En partant de z\'{e}ro pour concevoir un jeu complet --- architecture modulaire, graphismes th\'{e}matiques, syst\`{e}me de classes, et multijoueur en ligne --- nous avons collectivement franchi un cap significatif dans notre ma\^{i}trise de Python et de la conception logicielle.

Le r\'{e}sultat est un jeu jouable, coh\'{e}rent visuellement et narrativement, avec une rejouabilit\'{e} port\'{e}e par les six univers distincts et les cinq classes. Le module r\'{e}seau UDP, d\'{e}velopp\'{e} sans framework, constitue techniquement la partie la plus ambitieuse du projet.

Ce que ce projet nous a appris d\'{e}passe le code : planification r\'{e}aliste, communication r\'{e}guli\`{e}re, gestion des impr\'{e}vus, et fiert\'{e} collective d'un livrable abouti. Ces comp\'{e}tences nous accompagneront bien au-del\`{a} de la promo 2027.

\vspace{1cm}
\begin{center}
  \textcolor{mobiuspurple}{\large\bfseries M3G\_STUDIO --- MOBIUS Roguelike --- EPITA 2025}\\[4pt]
  \textcolor{mobiuscyan}{\textit{``Every loop ends. Until it doesn't.''}}
\end{center}

% ============================================================
%  ANNEXES
% ============================================================
\appendix
\chapter{Annexes illustr\'{e}es}

\textit{%
  Cette section est r\'{e}serv\'{e}e aux captures d'\'{e}cran et illustrations du jeu.
  Pour int\'{e}grer une image, remplacer le bloc \texttt{\textbackslash fbox} par :\\
  \texttt{\textbackslash includegraphics[width=0.8\textbackslash textwidth]\{screenshots/nom\_image.png\}}
}

\section{Menus}

\begin{figure}[H]
  \centering
  % \includegraphics[width=0.8\textwidth]{screenshots/menu_principal.png}
  \fbox{\parbox{0.8\textwidth}{\centering\vspace{3cm}
    \textit{[Illustration A1 : \'{E}cran titre MOBIUS]}
  \vspace{3cm}}}
  \caption{Annexe A1 --- Menu principal}
\end{figure}

\begin{figure}[H]
  \centering
  % \includegraphics[width=0.8\textwidth]{screenshots/mode_select.png}
  \fbox{\parbox{0.8\textwidth}{\centering\vspace{3cm}
    \textit{[Illustration A2 : S\'{e}lection de mode Solo / En ligne]}
  \vspace{3cm}}}
  \caption{Annexe A2 --- S\'{e}lection de mode}
\end{figure}

\begin{figure}[H]
  \centering
  % \includegraphics[width=0.8\textwidth]{screenshots/class_select.png}
  \fbox{\parbox{0.8\textwidth}{\centering\vspace{3cm}
    \textit{[Illustration A3 : S\'{e}lection de classe (5 cartes)]}
  \vspace{3cm}}}
  \caption{Annexe A3 --- S\'{e}lection de classe}
\end{figure}

\section{Gameplay --- \'{E}poques}

\begin{figure}[H]
  \centering
  % \includegraphics[width=0.8\textwidth]{screenshots/prehistoire.png}
  \fbox{\parbox{0.8\textwidth}{\centering\vspace{3cm}
    \textit{[Illustration B1 : Gameplay --- \`{E}re Pr\'{e}historique]}
  \vspace{3cm}}}
  \caption{Annexe B1 --- \`{E}re Pr\'{e}historique}
\end{figure}

\begin{figure}[H]
  \centering
  % \includegraphics[width=0.8\textwidth]{screenshots/edo.png}
  \fbox{\parbox{0.8\textwidth}{\centering\vspace{3cm}
    \textit{[Illustration B2 : Gameplay --- P\'{e}riode Edo (p\'{e}tales de cerisier)]}
  \vspace{3cm}}}
  \caption{Annexe B2 --- P\'{e}riode Edo}
\end{figure}

\begin{figure}[H]
  \centering
  % \includegraphics[width=0.8\textwidth]{screenshots/futuristique.png}
  \fbox{\parbox{0.8\textwidth}{\centering\vspace{3cm}
    \textit{[Illustration B3 : Gameplay --- \`{E}re Futuriste (scan holographique)]}
  \vspace{3cm}}}
  \caption{Annexe B3 --- \`{E}re Futuriste}
\end{figure}

\begin{figure}[H]
  \centering
  % \includegraphics[width=0.8\textwidth]{screenshots/ww2.png}
  \fbox{\parbox{0.8\textwidth}{\centering\vspace{3cm}
    \textit{[Illustration B4 : Gameplay --- Guerre Mondiale (pluie)]}
  \vspace{3cm}}}
  \caption{Annexe B4 --- Guerre Mondiale}
\end{figure}

\section{Multijoueur}

\begin{figure}[H]
  \centering
  % \includegraphics[width=0.8\textwidth]{screenshots/waiting_p2.png}
  \fbox{\parbox{0.8\textwidth}{\centering\vspace{3cm}
    \textit{[Illustration C1 : \'{E}cran En attente de P2 (host)]}
  \vspace{3cm}}}
  \caption{Annexe C1 --- Attente multijoueur (host)}
\end{figure}

\begin{figure}[H]
  \centering
  % \includegraphics[width=0.8\textwidth]{screenshots/coop.png}
  \fbox{\parbox{0.8\textwidth}{\centering\vspace{3cm}
    \textit{[Illustration C2 : Gameplay co-op 2 joueurs en ligne]}
  \vspace{3cm}}}
  \caption{Annexe C2 --- Partie co-op en ligne}
\end{figure}

\section{Communication \& planification}

\begin{figure}[H]
  \centering
  % \includegraphics[width=0.8\textwidth]{screenshots/site_web_full.png}
  \fbox{\parbox{0.8\textwidth}{\centering\vspace{3cm}
    \textit{[Illustration D1 : Site web M3G\_STUDIO]}
  \vspace{3cm}}}
  \caption{Annexe D1 --- Site web M3G\_STUDIO}
\end{figure}

\begin{figure}[H]
  \centering
  % \includegraphics[width=0.8\textwidth]{screenshots/gantt.png}
  \fbox{\parbox{0.8\textwidth}{\centering\vspace{3cm}
    \textit{[Illustration D2 : Diagramme de Gantt du projet]}
  \vspace{3cm}}}
  \caption{Annexe D2 --- Planification (Diagramme de Gantt)}
\end{figure}

\end{document}
